(1) 对状态方程在固定 T 下求导：
\[
P_v=-\frac{RT}{(v-b)^2}+\frac{2a}{v^3},\quad
P_{vv}=\frac{2RT}{(v-b)^3}-\frac{6a}{v^4}.
\]
在临界点，第一式给 \(RT_c=2a(v_c-b)^2/v_c^3\)，第二式给 \(RT_c=3a(v_c-b)^3/v_c^4\)。相除得 \(2v_c=3(v_c-b)\)，故 \(v_c=3b\)。回代得
\[
T_c=\frac{8a}{27Rb},\qquad P_c=\frac{a}{27b^2}.
\]

(2) 临界压缩因子
\[
Z_c=\frac{P_cv_c}{RT_c}
=\frac{(a/(27b^2))3b}{8a/(27b)}=\frac38.
\]
理想气体总满足 Z=1，说明排除体积与吸引共同使实际临界压缩系数偏离理想值。此数值是模型预测，不是任意实际流体的严格普适常数。

(3) 在单相稳定等温平衡区必须有 \((\partial P/\partial v)_T<0\)，等价于 \(\kappa_T=-1/(vP_v)>0\)。当 \(P_v>0\) 时形式上的 \(\kappa_T<0\) 表示均匀状态受机械扰动不稳定；当 \(P_v=0\) 压缩系数发散，需要考虑亚稳态、相分离与非线性效应，不能将整个 S 形等温线都当作稳定平衡。

(4) 在 T=Tc 且 v=3b+δv 附近，常数项为 Pc，线性及二次项因临界条件为零。三阶导数
\[
P_{vvv}=-\frac{6RT}{(v-b)^4}+\frac{24a}{v^5},
\quad
P_{vvv}(T_c,3b)=-\frac{a}{81b^5}.
\]
故
\[
P(T_c,v)=P_c-\frac{a}{486b^5}(\delta v)^3+O((\delta v)^4).
\]
这意味着临界等温线上 \(P-P_c\propto-(v-v_c)^3\)，体现此平均场模型的临界幂指数 δ=3，并与体积扰动的三次响应一致。
